\documentclass[11pt,a4paper]{article}

\usepackage[T1]{fontenc}
\usepackage[utf8]{inputenc}
\usepackage[provide=*,french]{babel}
\usepackage{amsmath,amssymb,amsfonts,mathtools}
\usepackage{booktabs,longtable,array,tabularx}
\usepackage{enumitem}
\usepackage{xcolor}
\usepackage{geometry}
\usepackage[expansion=false]{microtype}
\usepackage{hyperref}
\geometry{margin=2.25cm}
\hypersetup{
  colorlinks=true,
  linkcolor=blue!55!black,
  citecolor=green!40!black,
  urlcolor=blue!65!black,
  pdftitle={Transformers, régression in-context et problèmes inverses en EDP : revue critique},
  pdfauthor={Janis Aiad -- note de travail}
}

\newcommand{\R}{\mathbb{R}}
\newcommand{\E}{\mathbb{E}}
\newcommand{\cA}{\mathcal{A}}
\newcommand{\cC}{\mathcal{C}}
\newcommand{\cL}{\mathcal{L}}
\newcommand{\cT}{\mathcal{T}}
\newcommand{\cH}{\mathcal{H}}
\newcommand{\cX}{\mathcal{X}}
\newcommand{\cY}{\mathcal{Y}}
\newcommand{\op}{\mathrm{op}}
\newcommand{\tr}{\operatorname{tr}}
\newcommand{\rank}{\operatorname{rank}}
\newcommand{\cond}{\operatorname{cond}}
\newcommand{\ICL}{\mathrm{ICL}}
\newcommand{\HB}{\mathrm{HB}}
\newcommand{\PCG}{\mathrm{PCG}}
\newcommand{\FEM}{\mathrm{FEM}}
\newcommand{\OOD}{\mathrm{OOD}}
\newcommand{\TODO}{\textcolor{red!70!black}{\textbf{À vérifier}}}

\title{Transformers pour la régression in-context, les systèmes linéaires et les problèmes inverses en EDP\\
\large Revue critique, cartographie des lacunes et formulation de la question de recherche}
\author{Note de travail pour Janis Aiad}
\date{État de la littérature au 31 juillet 2026}

\begin{document}
\sloppy
\maketitle

\begin{abstract}
Cette note cartographie la littérature reliant l'apprentissage en contexte
(in-context learning, ICL), la régression linéaire, les Transformers vus comme
algorithmes d'optimisation, l'apprentissage d'opérateurs et les problèmes
directs ou inverses gouvernés par des équations aux dérivées partielles (EDP).
Elle distingue quatre objets souvent confondus : (i) apprendre directement un
surrogate de l'opérateur solution, (ii) identifier en contexte une loi ou un
opérateur caché, (iii) dérouler dans les activations un solveur numérique, et
(iv) apprendre un préconditionneur utilisé par un solveur externe. Le papier
de Takanami, Takahashi et Kabashima sur les tâches de régression de faible rang
est analysé en détail : son calcul de répliques prédit la loi du prédicteur et
une transition de phase en grande dimension, mais ne constitue ni une analyse
d'un solveur profond, ni un résultat d'EDP. La lacune centrale identifiée est
l'absence d'un cadre qui combine identification en contexte d'un opérateur
elliptique, décodeur itératif accéléré à poids partagés, préconditionnement
conditionné par le prompt, décomposition complète des erreurs et théorie
haute dimension du risque. Nous formulons en conséquence une question de
recherche et un protocole expérimental falsifiable autour d'un décodeur
HeavyBall/Chebyshev, avec PCG comme référence numérique et Richardson comme
baseline mécanistique.
\end{abstract}

\tableofcontents
\clearpage

\section{Périmètre, méthode et vocabulaire}

\subsection{Périmètre de la revue}

La revue couvre les travaux primaires disponibles jusqu'au 31 juillet 2026
dans les catégories suivantes :
\begin{enumerate}[leftmargin=*]
  \item ICL de fonctions linéaires, affines, parcimonieuses, de mélanges et de
  noyaux ;
  \item Transformers comme implémentations ou approximations d'algorithmes
  d'optimisation et d'algèbre linéaire ;
  \item profondeur, partage des poids, modèles bouclés et calcul au temps
  d'inférence ;
  \item théorie statistique non asymptotique, limites proportionnelles,
  matrices aléatoires, répliques et changement de distribution ;
  \item neural operators et Transformers pour EDP, y compris modèles de
  fondation et in-context operator learning ;
  \item problèmes inverses en EDP, préconditionneurs appris, multigrille appris
  et couplage réseaux--Krylov.
\end{enumerate}

Une exhaustivité littérale de tous les articles contenant le mot
« Transformer » serait peu informative. Le critère d'inclusion retenu est la
présence d'au moins un lien conceptuel direct avec l'un des objets suivants :
identification d'une tâche à partir d'exemples, résolution d'un système
linéaire dépendant de la tâche, apprentissage d'un opérateur entre espaces de
fonctions, inversion d'un paramètre d'EDP, ou accélération d'un solveur.

\subsection{Quatre sens différents de « Transformer pour résoudre une EDP »}

La première clarification est taxonomique. Les quatre paradigmes ci-dessous
ne répondent pas à la même question.

\begin{longtable}{p{0.18\textwidth}p{0.26\textwidth}p{0.24\textwidth}p{0.19\textwidth}}
\toprule
Paradigme & Entrée--sortie & Ce qui est appris & Garantie pertinente \\
\midrule
\endfirsthead
\toprule
Paradigme & Entrée--sortie & Ce qui est appris & Garantie pertinente \\
\midrule
\endhead
Surrogate d'opérateur & coefficient/champ initial $\mapsto$ solution & une
approximation globale de l'opérateur solution & approximation, généralisation,
stabilité et invariance à la discrétisation \\
ICL d'opérateur & prompt de paires fonctionnelles + requête $\mapsto$ réponse &
adaptation à une loi cachée sans mise à jour des poids & erreur en contexte,
diversité des tâches, transfert ID/OOD \\
Transformer-solveur & matrice ou prompt + second membre $\mapsto$ itérés & un
algorithme de résolution dans les activations & erreur par itération, stabilité
spectrale, coût par produit matrice--vecteur \\
Préconditionneur appris & opérateur/géométrie $\mapsto B_\theta$, puis solveur
externe & une transformation spectrale utile & SPD ou flexibilité du solveur,
conditionnement effectif, temps total \\
\bottomrule
\end{longtable}

DeepONet et FNO appartiennent d'abord à la première classe
\cite{lu2021deeponet,li2021fno,kovachki2023neuraloperator}. ICON appartient à la
deuxième \cite{yang2023icon,yang2024icongeneralization}. Les constructions de
gradient descent, Newton ou Richardson appartiennent à la troisième
\cite{vonoswald2023,fu2024secondorder,yan2026richardson}. NeuralIF et FCG-NO
appartiennent à la quatrième
\cite{hausner2023neuralif,rudikov2024fcgno}. Une contribution doit annoncer
explicitement sa classe, sinon une comparaison de performances ou de garanties
devient trompeuse.

\subsection{Deux inversions à ne pas confondre}

Pour une famille elliptique
\begin{equation}
  \cL_{a,V}u=-\nabla\!\cdot(a\nabla u)+Vu=f,
  \qquad u|_{\partial\Omega}=0,
  \label{eq:pde}
\end{equation}
deux problèmes inverses distincts apparaissent :
\begin{enumerate}[leftmargin=*]
  \item \emph{identification de loi} : retrouver $a$, $V$, ou une matrice de
  Galerkin $A$ à partir de paires $(f_i,u_i)$ ;
  \item \emph{inversion numérique} : pour un opérateur estimé $\widehat A$ et
  un nouveau second membre $f_\star$, calculer $\widehat A^{-1}f_\star$.
\end{enumerate}
Le premier est statistique et potentiellement mal posé ; le second est un
problème de solveur linéaire. Le projet considéré ici comporte les deux. Une
borne finale doit donc séparer l'erreur d'encodeur, l'erreur de discrétisation,
le biais de régularisation et l'erreur de profondeur finie du décodeur.

\section{ICL général : origine et grandes familles d'explications}

Brown et al. ont popularisé le comportement few-shot des grands modèles de
langage : une tâche est spécifiée par des exemples placés dans le prompt, sans
mise à jour des poids à l'inférence \cite{brown2020}. L'ICL ne se réduit
cependant pas à une seule cause. La littérature générale distingue au moins
quatre explications compatibles plutôt que mutuellement exclusives.

\paragraph{Inférence bayésienne implicite.}
Xie et al. modélisent le pré-entraînement comme l'apprentissage d'un mélange de
concepts latents ; conditionner sur un prompt revient alors à inférer le concept
courant \cite{xie2022bayesian}. Cette lecture est naturelle pour une famille
d'EDP : le prompt donne de l'information sur l'opérateur latent. Elle ne dit
pas encore comment inverser efficacement l'opérateur une fois identifié.

\paragraph{Mesa-optimisation.}
Dai et al. et von Oswald et al. interprètent la passe avant comme une procédure
d'optimisation interne \cite{dai2023meta,vonoswald2023}. Cette famille
d'explications motive directement un décodeur itératif. Elle doit être testée
contre des alternatives entrée--sortie équivalentes, car la proximité des
prédictions ne suffit pas à identifier un circuit.

\paragraph{Induction et structure de la distribution.}
Olsson et al. relient l'émergence de l'ICL à des têtes d'induction capables de
copier ou prolonger des motifs \cite{olsson2022}. Chan et al. montrent que des
propriétés de la distribution de séquences, notamment répétition et burstiness,
favorisent l'émergence de l'ICL \cite{chan2022distribution}. Un autre travail
de Chan et al. distingue la généralisation issue des informations stockées
dans les poids de celle issue du contexte \cite{chan2022weightscontext}. Pour
les EDP, l'analogue est la distinction entre prior de famille appris dans les
poids et opérateur particulier révélé par le prompt.

\paragraph{Rôle sémantique des démonstrations et apprenabilité.}
Min et al. montrent, sur des tâches de langage, que la distribution des entrées,
le format et l'espace des labels peuvent compter davantage que la correction
de chaque label individuel \cite{min2022demonstrations}. Wies et al. étudient
formellement l'apprenabilité de règles ICL par pré-entraînement
\cite{wies2023learnability}. Ces résultats invitent à ne pas supposer qu'un
prompt de paires $(f_i,u_i)$ est informatif par sa seule longueur : excitation
spectrale, diversité et format des tokens sont des variables de premier ordre.

Cette littérature générale fournit le contexte conceptuel, mais la régression
linéaire reste le laboratoire où mécanisme, entraînement et risque sont le plus
précisément calculables.

\section{ICL pour la régression linéaire : fondations empiriques et algorithmiques}

\subsection{Le paradigme expérimental fondateur}

Garg et al. ont établi le banc d'essai moderne : un Transformer est pré-entraîné
sur des prompts $(x_i,y_i)$ engendrés par des fonctions nouvelles à chaque
tâche, puis évalué sans mise à jour de paramètres sur une fonction non vue
\cite{garg2022}. Pour la régression linéaire gaussienne, des Transformers
standards atteignent empiriquement une performance proche des moindres carrés ;
ils peuvent aussi traiter régression parcimonieuse, petits réseaux et arbres de
décision. Ce travail établit une capacité, mais n'identifie pas de manière
unique l'algorithme interne.

Akyürek et al. ont ensuite donné trois niveaux de preuve
\cite{akyurek2023}: constructions explicites de gradient descent et de ridge en
forme fermée, correspondances empiriques entre prédictions et estimateurs
classiques, et décodage d'états internes ressemblant à des poids ou à des
moments. La leçon importante est que « le Transformer apprend la régression »
n'implique pas nécessairement « chaque couche est un pas de gradient » : le
même comportement entrée--sortie peut être réalisé par plusieurs circuits.

\subsection{Attention linéaire comme pas de gradient préconditionné}

Von Oswald et al. donnent une construction où une couche d'attention linéaire
réalise la transformation induite par un pas de gradient sur une perte de
régression \cite{vonoswald2023}. Ahn et al. montrent que les paramètres appris
peuvent correspondre à un gradient préconditionné dépendant de la distribution
des covariables \cite{ahn2023preconditioned}. Mahankali et al. prouvent que,
pour une couche d'auto-attention linéaire et des covariables gaussiennes, le
minimiseur populationnel implémente précisément un pas de gradient, ou un pas
préconditionné en cas d'anisotropie \cite{mahankali2023}.

Ces résultats sont centraux pour Richardson, car le pas
\[
  z_{\ell+1}=z_\ell+B(c-Hz_\ell)
\]
est exactement un pas de gradient préconditionné pour une quadratique. Ils ne
prouvent cependant pas qu'un Transformer standard entraîné sur des EDP apprend
spontanément le bon $B$, ni que plusieurs couches reproduisent un solveur
stable hors distribution.

\subsection{Ce que l'entraînement apprend réellement}

Zhang, Frei et Bartlett analysent le gradient flow d'une couche linéaire et
montrent convergence vers un minimum global et robustesse à certains shifts,
mais fragilité aux shifts de covariance des covariables
\cite{zhang2023trained}. Les travaux de Wu et al. quantifient le nombre de
tâches de pré-entraînement requis pour la régression linéaire
\cite{wu2024tasks}. Raventós et al. montrent empiriquement et théoriquement que
la diversité des tâches gouverne une transition entre mémorisation de tâches et
algorithme non bayésien généralisant à de nouvelles tâches
\cite{raventos2023}. Li et al. étudient stabilité et généralisation de
Transformers vus comme algorithmes \cite{li2023algorithms}.

Les résultats de sélection d'algorithme de Bai et al. sont complémentaires :
des Transformers suffisamment construits peuvent sélectionner en contexte un
algorithme adapté parmi plusieurs classes statistiques
\cite{bai2023statisticians}. Ils portent surtout sur expressivité et garanties
de risque, pas sur l'identification mécanistique d'un modèle effectivement
entraîné.

\subsection{Au-delà du premier ordre}

La conclusion « les Transformers font seulement gradient descent » est trop
forte. Fu et al. observent des taux proches d'une itération de Newton et donnent
une construction de plusieurs itérations de Newton en profondeur
\cite{fu2024secondorder}. Vladymyrov et al. montrent la polyvalence de
l'attention linéaire comme apprenant en contexte \cite{vladymyrov2024}. Pour
les fonctions non linéaires, des constructions relient attention et gradient
fonctionnel \cite{cheng2024functional}; pour la régression par noyau gaussien,
Yan et al. donnent en 2026 une construction Transformer softmax réalisant une
itération de Richardson préconditionnée \cite{yan2026richardson}.

Le papier de Yan et al. est particulièrement proche du projet, mais sa portée
doit être lue précisément. L'attention softmax produit un opérateur de noyau
gaussien normalisé ; des MLP ReLU approximent des opérations scalaires locales
(produits, inverse et normalisations). La construction utilise
$O(\log(1/\varepsilon))$ blocs et une largeur MLP dépendant de la précision.
L'alignement empirique couche--itéré est le plus fort pour Richardson. Cela
justifie Richardson comme mécanisme réalisable, mais ne montre pas qu'il est le
meilleur solveur à budget de calcul fixé.

\section{Profondeur, boucles et conditionnement}

\subsection{Pourquoi une couche n'est pas une itération universelle}

Les constructions d'expressivité fixent généralement des paramètres à la main.
Elles prouvent qu'un circuit existe, non qu'il est identifiable par
l'entraînement. De plus, les couches peuvent effectuer des sous-opérations
différentes : lecture du prompt, agrégation de moments, mise à jour d'état et
readout. Compter « une couche = une itération » n'est honnête que si le bloc et
l'état ont été définis pour rendre l'égalité exacte.

Le conditionnement spectral rend cette nuance essentielle. Sur une quadratique
SPD de condition $\kappa$, Richardson optimal demande typiquement
$O(\kappa\log(1/\varepsilon))$ itérations, tandis que Chebyshev, HeavyBall
optimal et CG présentent une dépendance en $O(\sqrt\kappa
\log(1/\varepsilon))$ sous leurs hypothèses classiques
\cite{polyak1964,saad2003}. Une architecture qui apprend seulement un pas de
gradient peut donc être mécanistiquement correcte mais numériquement dominée.

\subsection{Transformers bouclés et poids partagés}

Gatmiry et al. analysent des Transformers linéaires bouclés : le minimum global
de la perte populationnelle réalise un gradient préconditionné multi-pas, et le
partage des poids impose un biais vers les itérations de point fixe
\cite{gatmiry2024learn}. Leur étude complémentaire montre que profondeur et
diversité des tâches interagissent avec le conditionnement, et que les modèles
bouclés ont de meilleures propriétés de robustesse et de monotonie en
profondeur que certains modèles non partagés \cite{gatmiry2024depth}.
Chen et al. obtiennent également des gains d'efficacité théorique pour des
Transformers bouclés réalisant plusieurs pas de gradient
\cite{chen2024looped}. Huang et al. utilisent une chaîne de pensée explicite
pour apprendre plusieurs pas de gradient \cite{huang2025cot}.

Pour le projet, le partage des poids est plus qu'une économie de paramètres :
il rend plausible une extrapolation en profondeur et permet d'énoncer une loi
de récurrence indépendante de $L$. HeavyBall à coefficients constants s'insère
naturellement dans ce cadre. Chebyshev nécessite des coefficients dépendant de
l'itération, donc soit des couches non partagées, soit un état de profondeur ou
un contrôleur de coefficients.

\subsection{Régression structurée, parcimonieuse et mélanges}

La régression linéaire isotrope n'épuise pas les structures utiles en EDP.
Chen, Zhao et Zou montrent que plusieurs têtes peuvent servir à prétraiter une
régression parcimonieuse avant des pas d'optimisation plus simples
\cite{chen2024sparse}. Jin et al. analysent les mélanges de régressions, où le
Transformer doit implicitement identifier une composante avant d'estimer ses
coefficients \cite{jin2024mixture}. Ces modèles sont proches d'une famille
d'opérateurs à plusieurs régimes. Les travaux sur les représentations et les
sous-espaces faibles dimensions expliquent aussi comment un ICL peut exploiter
des caractéristiques partagées plutôt qu'une covariance isotrope
\cite{guo2023representations,kwon2025subspace}.

\section{Théorie haute dimension : matrices aléatoires et répliques}

\subsection{Limite proportionnelle de l'attention linéaire}

Lu et al. développent une théorie asymptotique exactement soluble pour une
couche d'attention linéaire en régression \cite{lu2025asymptotic}. Dimension
des tokens, longueur de contexte et diversité des tâches croissent
proportionnellement, tandis que le nombre d'exemples de pré-entraînement suit
une autre échelle. La théorie prédit double descente et transition entre
mémorisation et véritable ICL. Letey et al. étendent ce point de vue au
décalage entre covariance de tâches de pré-entraînement et de test, et
introduisent une mesure d'alignement qui prédit la généralisation
\cite{letey2026alignment}.

Ce régime complète les bornes non asymptotiques : celles-ci contrôlent le pire
cas ou une probabilité haute, tandis que la limite proportionnelle vise une
valeur précise du risque typique. Les deux niveaux ne sont pas interchangeables.

\subsection{Le papier de faible rang 2510.04548}

Takanami, Takahashi et Kabashima étudient une couche d'attention linéaire
entraînée sur des tâches
\[
  w^\mu=\sqrt{D/r}\,A v^\mu,\qquad A\in\R^{D\times r},
\]
partageant un sous-espace latent de rang $r$ \cite{takanami2026}. Leur limite
fait croître $D$, la longueur de contexte, le nombre de tâches de base, le rang
et le nombre d'instances en gardant fixes plusieurs rapports :
\[
 \alpha=L/D,\qquad \rho=r/D,\qquad
 \kappa=M_0/D,\qquad \gamma=M/(DM_0).
\]
Le calcul de répliques prédit une décomposition en loi
\[
 \widehat y\overset{d}{=}
 \widehat y_{\mathrm{algo}}+
 \widehat y_{\mathrm{mem}}+
 \widehat y_{\mathrm{struct}}.
\]
Le terme algorithmique projette un estimateur de type matched filter sur le
sous-espace appris $AA^\top$ ; les deux autres termes représentent un bruit lié
aux tâches mémorisées et un bruit lié au défaut d'alignement avec la structure.
Le papier prédit aussi :
\begin{itemize}[leftmargin=*]
  \item une transition lorsque la diversité de tâches devient suffisante par
  rapport au rang latent ;
  \item un compromis spécialisation--robustesse OOD ;
  \item une régularisation implicite d'ordre $1/\alpha$ créée par les
  fluctuations d'un contexte de pré-entraînement fini ;
  \item un accord numérique des paramètres d'ordre et du risque avec les
  expériences, plus une vérification qualitative sur attention softmax.
\end{itemize}

Il faut toutefois employer la terminologie des auteurs : il s'agit de
\emph{Results} issus de la méthode des répliques, et non de théorèmes rigoureux,
car l'argument utilise continuation analytique, échange de limites et ansatz de
symétrie des répliques. Le résultat est extrêmement pertinent pour modéliser la
phase d'encodeur ou l'apprentissage d'un sous-espace de préconditionnement. Il
ne traite pas :
\begin{enumerate}[leftmargin=*]
  \item une attention profonde ou bouclée ;
  \item une matrice de Galerkin ou un opérateur SPD dépendant de la tâche ;
  \item une propagation d'erreur de $\widehat A$ vers $\widehat A^{-1}f$ ;
  \item un polynôme de solveur Richardson, HeavyBall, Chebyshev ou CG ;
  \item le coût ou la stabilité d'un préconditionneur appris.
\end{enumerate}

La bonne utilisation dans le projet est donc de reprendre sa machinerie pour
la loi de l'encodeur ou du préconditionneur, puis de composer cette loi avec un
calcul spectral exact du décodeur. Le risque du décodeur fixe peut souvent
s'écrire comme une intégrale spectrale pondérée ; c'est ce chaînon qui manque
au papier de faible rang.

\subsection{Rigueur possible au-delà des répliques}

La méthode des répliques est un calcul heuristique puissant, mais certaines de
ses prédictions peuvent être rendues rigoureuses par interpolation adaptative
dans des modèles bayésiens ou par état d'évolution AMP
\cite{barbier2019adaptive,bayati2011amp}. Pour un projet de solveur, trois
niveaux de résultat peuvent coexister proprement :
\begin{enumerate}[leftmargin=*]
  \item un théorème déterministe conditionnel en fonction du spectre de
  $B_\theta\widehat A$ ;
  \item une borne probabiliste non asymptotique pour l'encodeur ;
  \item une prédiction replica/RMT précise du risque typique en limite
  proportionnelle, explicitement annoncée comme telle.
\end{enumerate}
Cette séparation évite de transformer une formule de répliques en garantie
uniforme injustifiée.

\section{Des systèmes linéaires aux opérateurs et aux EDP}

\subsection{ICL de systèmes linéaires}

Cole, Lu, Xu et Zhang étudient directement l'ICL de systèmes linéaires avec
une couche de Transformer linéaire \cite{cole2025systems}. Ils donnent des lois
d'échelle en fonction du nombre de tâches de pré-entraînement $N$, de la
longueur d'entraînement $n$ et de la longueur de prompt test $m$, ainsi qu'une
condition de diversité nécessaire et suffisante dans leur cadre pour certains
shifts de tâches. Leur théorème abstrait transfère le résultat fini-dimensionnel
à l'apprentissage d'opérateurs et donne une application aux EDP elliptiques
linéaires après projection de Galerkin.

Ce travail est la référence théorique la plus proche, mais le modèle est un
prédicteur linéaire à une couche, pas un solveur récurrent explicitement
analysé. Sa borne regroupe l'effet statistique et l'approximation de
discrétisation ; elle ne compare pas les complexités en profondeur de
Richardson, HeavyBall, Chebyshev et PCG, et n'apprend pas explicitement un
préconditionneur SPD conditionné par le prompt.

Le prolongement aux systèmes dynamiques montre par ailleurs une séparation de
profondeur : des Transformers multi-couches atteignent l'échelle d'un
estimateur moindres carrés alors qu'une classe à une couche conserve une erreur
non nulle sur données non-i.i.d. \cite{cole2025dynamics}. Des travaux voisins
interprètent les Transformers comme estimateurs d'état implicites ou étudient
la représentation du filtre de Kalman
\cite{goel2024kalman,akram2024state}. Ils confirment qu'un contexte corrélé
change qualitativement le problème par rapport à la régression i.i.d.

\subsection{Operator learning classique}

DeepONet sépare un réseau \emph{branch}, qui encode la fonction d'entrée
échantillonnée, et un réseau \emph{trunk}, qui encode le point de sortie
\cite{lu2021deeponet}. FNO paramètre l'opérateur intégral dans l'espace de
Fourier et vise une certaine invariance à la résolution
\cite{li2021fno}. Le cadre neural operator général formalise des couches entre
espaces de fonctions et leurs propriétés d'approximation
\cite{kovachki2023neuraloperator}. Des estimations d'erreur pour DeepONet
séparent encodage, approximation et reconstruction
\cite{lanthaler2022error}; des bornes de complexité de Rademacher ont aussi été
développées pour FNO \cite{kim2022rademacher}.

Ces travaux apprennent habituellement une loi amortie sur une distribution
d'opérateurs ou de paramètres. Ils ne réalisent pas nécessairement une
adaptation ICL à un opérateur entièrement nouveau à partir de quelques paires
$(f_i,u_i)$.

\subsection{Transformers comme neural operators}

Le Galerkin Transformer montre que l'attention linéarisée sans softmax peut être
interprétée comme une projection de type Petrov--Galerkin et l'applique à des
problèmes directs et à une identification de coefficient d'interface
\cite{cao2021galerkin}. OFormer développe une architecture Transformer pour les
opérateurs spatio-temporels \cite{li2022oformer}. GNOT traite entrées
hétérogènes, maillages irréguliers et décomposition géométrique souple
\cite{hao2023gnot}. Transolver construit des tokens physiques et une attention
efficace pour géométries générales \cite{wu2024transolver}.

Ces architectures sont des surrogates puissants ; leur attention n'implique pas
à elle seule une itération de Galerkin ou de Krylov. La ressemblance entre
attention et projection variationnelle est structurelle, pas une preuve de
convergence vers la solution d'un système donné.

\subsection{Pré-entraînement multi-EDP et modèles de fondation}

DPOT combine Transformer d'opérateur, attention de Fourier et pré-entraînement
auto-régressif débruitant sur de nombreuses trajectoires
\cite{hao2024dpot}. Poseidon utilise un Transformer d'opérateur multi-échelle
et la propriété de semi-groupe pour pré-entraînement et transfert sur plusieurs
EDP \cite{herde2024poseidon}. PDEformer représente symboliquement l'EDP comme
un graphe, traité par Transformer, puis décode une représentation implicite de
la solution ; il fournit aussi des expériences d'identification de coefficients
\cite{ye2024pdeformer}. UPS et Unisolver visent également l'adaptation
multi-physique et les changements de modalité
\cite{shen2024ups,lu2024unisolver}.

Leurs résultats portent surtout sur accuracy, sample efficiency, zero-shot ou
fine-tuning. Ils ne donnent généralement pas une décomposition complète
\[
  \text{discrétisation} + \text{identification de loi}
  + \text{solveur fini} + \text{méta-généralisation}.
\]
Ils constituent en revanche des baselines empiriques nécessaires si la
contribution revendique un avantage pratique général en EDP.

\subsection{In-context operator learning}

ICON reçoit des démonstrations fonctionnelles comme prompt et peut réaliser
différentes prédictions directes ou inverses sans fine-tuning
\cite{yang2023icon}. Une étude sur les lois de conservation 1D montre une
généralisation à certaines formes d'EDP non vues, mais aussi la dépendance au
design du prompt et au domaine de tâches \cite{yang2024icongeneralization}.
Cole et al. fournissent la première justification mathématique directe de ce
paradigme pour des opérateurs linéaires \cite{cole2025systems}. Mishra, Patel et
Tewari montrent enfin que des Continuum Transformers peuvent réaliser un
gradient d'opérateur en contexte \cite{mishra2025continuum}.

Cette branche est la plus proche du projet au niveau de la sémantique du
prompt. La lacune demeure le décodeur numérique : ICON apprend une règle de
prédiction, alors que le projet cherche à exposer et accélérer la résolution du
système conditionné par le prompt.

\section{Problèmes inverses en EDP et architectures Transformer}

\subsection{Surrogates inverses et opérateurs inverses}

Les problèmes inverses en EDP sont mal posés : stabilité conditionnelle,
régularisation et bruit d'observation ne peuvent être ignorés
\cite{engl1996inverse,stuart2010bayesian}. Un réseau qui prédit un coefficient
à partir d'une solution n'est donc pas seulement un solveur de système
linéaire ; il apprend un inverse régularisé dépendant du prior et du modèle de
bruit.

ViTO utilise une architecture Vision Transformer--Operator pour des problèmes
inverses paramétriques \cite{ovadia2023vito}. Les Neural Inverse Operators
apprennent directement une application données--paramètre pour des inversions
d'EDP \cite{molinaro2023nio}. BiLO imbrique un opérateur local dans une
optimisation bi-niveau pour l'identification de paramètres
\cite{cao2024bilo}. LNO réalise des calculs directs et inverses dans un espace
latent via cross-attention physique \cite{wang2024lno}. PDEformer inclut aussi
la récupération de coefficients \cite{ye2024pdeformer}.

Ces méthodes doivent être comparées sur l'identification de $a$ ou $V$, mais
elles ne répondent pas automatiquement à la question du meilleur décodeur de
$\widehat A u=f$. Inverse statistique et solveur numérique doivent rester deux
étages auditables.

\subsection{Formulation faible et identification de faible dimension}

Si
\[
  a_z(x)=a_0(x)+\sum_{k=1}^K z_k\psi_k(x),
\]
les paires de prompt $(f_i,u_i)$ et des fonctions test $v_r$ donnent
\[
  \sum_{k=1}^K z_k
  \int_\Omega \psi_k\nabla u_i\!\cdot\nabla v_r
  =
  \int_\Omega f_i v_r-
  \int_\Omega a_0\nabla u_i\!\cdot\nabla v_r.
\]
Après empilement, $G_mz\simeq b_m$. Cette réduction relie le problème inverse
elliptique à une régression linéaire structurée, mais les lignes de $G_m$ ne
sont ni isotropes ni indépendantes de la tâche. Leur covariance dépend de la
PDE, du choix de tests, de l'excitation des modes par les sources et du bruit
sur $u_i$. Les théories ICL gaussiennes ne s'appliquent donc pas par simple
substitution ; il faut contrôler ou modéliser ce design structuré.

\section{Préconditionnement appris et solveurs hybrides}

\subsection{Pourquoi cette littérature est une baseline obligatoire}

La littérature « learned solver » adopte souvent une stratégie plus sûre que
l'imitation complète d'un algorithme : le réseau produit un warm start, une
correction de grille grossière, une factorisation approchée ou un
préconditionneur ; un solveur classique impose ensuite la précision finale.
Cette séparation offre des résidus vérifiables et permet de conserver des
garanties numériques.

Sappl et al. apprennent des préconditionneurs pour CG
\cite{sappl2019preconditioners}. Li et al. apprennent un préconditionneur de
CG pour des systèmes issus d'EDP \cite{li2023preconditioner}. NeuralIF produit
par GNN une factorisation incomplète parcimonieuse et l'entraîne via un proxy
Frobenius efficace \cite{hausner2023neuralif}. FCG-NO utilise un neural
operator potentiellement non linéaire comme préconditionneur d'un flexible CG,
avec transfert entre résolutions \cite{rudikov2024fcgno}. UGrid apprend un
schéma multigrille avec contrôle de convergence \cite{han2024ugrid}. Les
travaux de multigrille appris remontent au moins à l'apprentissage de
prolongations et de coarse operators par Greenfeld et al.
\cite{greenfeld2019multigrid}.

Ces méthodes sont les vraies baselines d'une revendication « meilleur solveur
decoder ». FNO, DeepONet ou ICON seuls ne suffisent pas comme comparateurs : il
faut aussi mesurer contre Jacobi/ILU/AMG, CG/PCG, Chebyshev, warm start neural
et préconditionneur neural + Krylov.

\subsection{Trois architectures possibles pour le projet}

\paragraph{A. Transformer pur à polynôme fixe.}
Le bloc partagé conserve $(z_\ell,z_{\ell-1})$ et applique
\begin{equation}
z_{\ell+1}=z_\ell+\alpha B_\theta(\cC)
(c-Hz_\ell)+\beta(z_\ell-z_{\ell-1}).
\label{eq:hb}
\end{equation}
L'attention calcule le produit prompt--vecteur ou l'agrégation faible ; le
mélange linéaire des canaux réalise les coefficients scalaires. Aucune division
adaptative n'est requise. Avec $\beta=0$, on retrouve Richardson. Avec
$\beta>0$, on obtient HeavyBall. Des coefficients dépendant de $\ell$ donnent
une semi-itération de Chebyshev.

\paragraph{B. Transformer-conditioned PCG.}
Le Transformer produit un préconditionneur SPD, par exemple
\[
 B_\theta(\cC)=D^{-1/2}
 (I+U_\theta(\cC)U_\theta(\cC)^\top)D^{-1/2},
\]
puis une cellule PCG externe effectue exactement les produits scalaires,
quotients et conjugaisons. C'est probablement la meilleure architecture
pratique, mais elle doit être nommée \emph{Transformer-conditioned PCG} et non
« vanilla Transformer implements PCG ».

\paragraph{C. Transformer pur imitant PCG.}
Il faut stocker résidu, résidu préconditionné, direction, plusieurs produits
scalaires et des quotients dépendant du prompt et de l'itération. Une
construction par approximation MLP est possible sur un domaine dont les
dénominateurs sont séparés de zéro, mais elle ajoute une erreur arithmétique et
plusieurs sous-blocs par itération. Cette voie est théoriquement originale,
mais moins naturelle pour le premier démonstrateur.

\subsection{HeavyBall, Chebyshev et PCG : comparaison honnête}

Pour $S=B^{1/2}HB^{1/2}$ avec spectre dans $[\mu,L_s]$, HeavyBall constant
obéit mode par mode à
\[
 e_{\ell+1}(\lambda)=
 (1+\beta-\alpha\lambda)e_\ell(\lambda)
 -\beta e_{\ell-1}(\lambda).
\]
Les paramètres quadratiques classiques sont
\[
 \alpha_\star=\frac{4}{(\sqrt{L_s}+\sqrt\mu)^2},\qquad
 \beta_\star=
 \left(\frac{\sqrt{L_s}-\sqrt\mu}
 {\sqrt{L_s}+\sqrt\mu}\right)^2.
\]
Le taux asymptotique modal est
$q=(\sqrt\kappa-1)/(\sqrt\kappa+1)$, mais une borne finie doit conserver le
facteur transitoire lié aux racines doubles aux extrémités. Chebyshev minimise
un polynôme sur un intervalle spectral connu et constitue la meilleure
baseline polynomiale non adaptative. PCG choisit ses coefficients en fonction
du second membre et du sous-espace de Krylov ; il reste la baseline supérieure
pour SPD, mais sa non-linéarité en $f$ complique une théorie replica du décodeur.

Le meilleur compromis scientifique est donc : HeavyBall partagé comme
contribution « Transformer pur », Chebyshev comme contrôle polynomial, PCG
comme oracle numérique, et Transformer-conditioned PCG comme architecture
hybride orientée performance.

\section{Décomposition des erreurs : ce qu'une théorie complète doit prouver}

Soit $A$ l'opérateur discret vrai, $\widehat A$ l'opérateur encodé, et
$\widehat u_L$ la sortie après $L$ boucles. Une décomposition minimale est
\begin{align}
\E\|\widehat u_L-u\|_{\cH}^2
\lesssim{}&
\underbrace{\varepsilon_{\mathrm{disc}}(d)}_{\text{Galerkin/maillage}}
+\underbrace{\varepsilon_{\mathrm{inv}}(m,\sigma,K)}_{\text{identification de loi}}
\nonumber\\
&+\underbrace{\varepsilon_{\mathrm{meta}}(N,n)}_{\text{pré-entraînement}}
+\underbrace{\varepsilon_{\mathrm{solver}}(L,B_\theta\widehat A)}_{\text{profondeur finie}}
+\underbrace{\varepsilon_{\mathrm{arch}}}_{\text{réalisation par attention}}.
\label{eq:riskdecomp}
\end{align}

Pour un décodeur polynomial fixe, conditionnellement à $\widehat A$ et
$B_\theta$, la sortie est linéaire en $f$. Si $p_L$ est son polynôme résiduel,
\[
 \widehat u_L=widehat A^{-1}
 [I-p_L(B_\theta\widehat A)]f.
\]
Pour $f\sim\mathcal N(0,\Sigma_f)$, le risque conditionnel du solveur admet une
trace exacte
\[
 \E_f\|\widehat u_L-\widehat A^{-1}f\|_2^2
 =\tr(\Sigma_f E_L^\top E_L),
 \quad E_L=\widehat A^{-1}p_L(B_\theta\widehat A).
\]
Cette formule rend HeavyBall/Chebyshev compatibles avec une intégration RMT ou
replica de la loi spectrale de l'encodeur. PCG, dont les coefficients dépendent
de $f$, demande des paramètres d'ordre supplémentaires sur le processus de
Krylov.

La propagation opérateur--solution suit l'identité de résolvante
\[
 \widehat A^{-1}-A^{-1}
 =\widehat A^{-1}(A-\widehat A)A^{-1}.
\]
Sous coercivité uniforme $A,\widehat A\succeq a_0I$, une borne d'opérateur pour
$\widehat A-A$ devient une borne de solution avec un facteur $a_0^{-2}$. Ce
pont est indispensable : une bonne perte Frobenius sur l'opérateur n'est pas
automatiquement une bonne erreur de solution si les erreurs sont alignées avec
des modes spectraux fragiles.

\section{Tableau critique de la littérature la plus proche}

\begin{longtable}{p{0.19\textwidth}p{0.18\textwidth}p{0.22\textwidth}p{0.26\textwidth}}
\toprule
Référence & Objet & Résultat central & Ce qui manque pour notre question \\
\midrule
\endfirsthead
\toprule
Référence & Objet & Résultat central & Ce qui manque pour notre question \\
\midrule
\endhead
Garg et al. \cite{garg2022} & ICL régression & Transformer entraîné proche de
moindres carrés & mécanisme unique, EDP, solveur spectral \\
Akyürek et al. \cite{akyurek2023} & algorithmes ICL & constructions GD, ridge,
OLS & apprentissage d'un opérateur PDE et stabilité OOD \\
Ahn et al. \cite{ahn2023preconditioned} & régression anisotrope & gradient
préconditionné & profondeur accélérée et contrainte SPD \\
Fu et al. \cite{fu2024secondorder} & second ordre & construction et alignement
Newton & arithmétique, EDP et erreur encodeur--décodeur \\
Gatmiry et al. \cite{gatmiry2024learn} & Transformer bouclé & multi-pas GD
appris & acceleration HB/Chebyshev et problème inverse \\
Lu et al. \cite{lu2025asymptotic} & théorie haute dimension & risque exact,
double descente, diversité & opérateur SPD, profondeur et solveur \\
Takanami et al. \cite{takanami2026} & tâches de faible rang & décomposition
replica et transition de phase & décodeur profond et inversion $A^{-1}$ \\
Cole et al. \cite{cole2025systems} & systèmes linéaires/EDP & bornes ID/OOD et
transfert opérateur & solveur récurrent accéléré et préconditionneur appris \\
Yan et al. \cite{yan2026richardson} & KRR softmax & Richardson constructif avec
garantie & meilleur solveur, EDP et apprentissage SPD \\
ICON \cite{yang2023icon} & ICL d'opérateur & adaptation sans fine-tuning &
garantie de solveur et coût spectral \\
Poseidon/DPOT \cite{herde2024poseidon,hao2024dpot} & fondation EDP & transfert
empirique multi-physique & mécanisme de résolution et théorie de risque \\
NeuralIF \cite{hausner2023neuralif} & préconditionneur CG & factorisation
apprise efficace & ICL par prompt et théorie meta/OOD \\
FCG-NO \cite{rudikov2024fcgno} & neural operator + FCG & préconditionneur non
linéaire multi-résolution & Transformer pur et décomposition replica \\
\bottomrule
\end{longtable}

\section{La lacune précise et la question que nous améliorons}

\subsection{Lacune}

Aucune des branches examinées ne réunit simultanément les cinq propriétés
suivantes :
\begin{enumerate}[leftmargin=*]
  \item une tâche est un opérateur elliptique caché, identifié uniquement par
  un prompt de paires source--solution ;
  \item le modèle se généralise à un nouvel opérateur sans mise à jour de ses
  poids ;
  \item le décodeur est un solveur accéléré explicite, bouclé et à poids
  partagés, plutôt qu'un surrogate opaque ;
  \item le préconditionnement est appris à partir du prompt avec une propriété
  SPD ou une stabilité spectrale vérifiable ;
  \item la théorie sépare erreur de Galerkin, erreur d'encodeur, erreur de
  profondeur, réalisation architecturale et généralisation ID/OOD, avec en
  plus une prédiction haute dimension du risque typique.
\end{enumerate}

\subsection{Question principale proposée}

\begin{quote}
\textbf{Peut-on pré-entraîner un Transformer bouclé qui, à partir d'un prompt
de paires source--solution générées par un opérateur elliptique inconnu,
identifie les directions spectrales utiles et applique à un nouveau second
membre un décodeur préconditionné accéléré à poids partagés, de sorte que (i)
sa trajectoire interne réalise exactement ou uniformément HeavyBall/Chebyshev,
(ii) son erreur totale admet une décomposition non asymptotique
encodeur--solveur, (iii) son risque typique en grande dimension est prédit par
RMT/répliques, et (iv) il surpasse Richardson et les neural operators directs à
budget égal tout en approchant PCG ?}
\end{quote}

Cette formulation améliore la question de Cole et al. en remplaçant la seule
prédiction d'un système linéaire par une architecture encodeur--solveur
explicite et accélérée. Elle améliore la ligne von Oswald--Ahn--Gatmiry en
passant d'une régression synthétique à un design faible issu d'EDP et en
distinguant apprentissage du préconditionneur et itération. Elle améliore la
ligne Takanami--Lu en ajoutant une dynamique de décodeur profonde dont le
polynôme spectral est observable. Enfin, elle se distingue de FCG-NO en
imposant adaptation ICL à une loi cachée et en visant un circuit Transformer
pur avant l'hybride PCG.

\subsection{Sous-questions falsifiables}

\begin{enumerate}[leftmargin=*]
  \item \textbf{Réalisation.} Un bloc partagé sans MLP arithmétique peut-il
  réaliser l'agrégation $G_m^\top(b_m-G_mz_\ell)$ et la mémoire HeavyBall avec
  une erreur uniforme contrôlée ?
  \item \textbf{Apprentissage spectral.} Le préconditionneur appris réduit-il
  réellement $\cond(B_\theta^{1/2}H B_\theta^{1/2})$, ou seulement la perte sur
  la distribution de seconds membres vue ?
  \item \textbf{Structure de faible rang.} Le modèle retrouve-t-il un sous-espace
  de correction analogue au $AA^\top$ de Takanami et al., et existe-t-il une
  transition en fonction du rang latent, de la diversité et du nombre de
  prompts ?
  \item \textbf{Robustesse.} Les coefficients partagés restent-ils stables sous
  shift de maillage, covariance des sources, coefficients d'EDP et condition
  number ?
  \item \textbf{Performance.} À nombre égal de produits $Hv$, HeavyBall appris
  bat-il Richardson ; Chebyshev oracle borne-t-il le meilleur polynôme fixe ;
  combien reste-t-il d'écart à PCG, AMG et neural-preconditioned CG ?
  \item \textbf{Théorie.} La loi replica de l'encodeur, composée avec le
  polynôme du décodeur, prédit-elle le risque couche par couche et ses pics
  transitoires ?
\end{enumerate}

\section{Contributions réalistes et ordre de travail}

Une hiérarchie crédible de contributions est la suivante.

\paragraph{Contribution 1 : mécanisme exact.}
Construire un état de tokens contenant $(z_\ell,z_{\ell-1})$, des tokens de
prompt faibles et éventuellement des tokens spectraux. Prouver qu'un bloc
partagé implémente un pas de HeavyBall préconditionné, modulo l'erreur déjà
nécessaire pour le produit prompt--vecteur. Cette construction ne requiert ni
inverse scalaire adaptatif ni approximation de quotient.

\paragraph{Contribution 2 : transfert encodeur--décodeur.}
Établir un théorème conditionnel qui compose : coercivité et Galerkin,
concentration de $\widehat A$, stabilité de $B_\theta$, polynôme résiduel
$p_L$, puis risque de solution. Comparer la borne et la trace exacte à
Richardson.

\paragraph{Contribution 3 : théorie haute dimension.}
Adapter le modèle de faible rang aux lignes de design faible $G_m$, obtenir la
loi du sous-espace ou du préconditionneur appris, puis intégrer
$p_L(\lambda)^2$ contre la mesure spectrale limite. Cette partie peut être
annoncée comme prédiction replica, avec une validation numérique systématique.

\paragraph{Contribution 4 : benchmark solveur.}
Comparer, à même nombre de produits matrice--vecteur et même budget mémoire :
Richardson scalaire, Jacobi--Richardson, HeavyBall appris/oracle, Chebyshev,
CG/PCG, AMG/ILU, NeuralIF ou préconditionneur neuronal, FCG-NO, ainsi que des
surrogates directs FNO/GNOT/Poseidon lorsque le jeu de données le permet.

\paragraph{Contribution 5 : hybride performance.}
Une fois le mécanisme pur établi, entraîner un Transformer qui produit un
préconditionneur SPD pour un PCG externe. Cette version peut viser la meilleure
performance pratique sans diluer la contribution théorique du décodeur pur.

\section{Protocole expérimental minimal pour soutenir les revendications}

\subsection{Axes de données}

Les expériences doivent varier séparément : dimension de Galerkin $d$, rang
latent $K$, longueur du prompt $m$, nombre de tâches $N$, bruit, condition
number, covariance des sources, diversité des opérateurs, géométrie et
résolution. Les shifts ID/OOD doivent inclure :
\begin{itemize}[leftmargin=*]
  \item nouvelles réalisations du même prior de coefficients ;
  \item rang ou longueur de corrélation hors entraînement ;
  \item contraste elliptique plus élevé ;
  \item maillage ou base différents ;
  \item nouvelles covariances de seconds membres ;
  \item combinaisons de directions spectrales rares.
\end{itemize}

\subsection{Métriques}

La seule MSE de sortie est insuffisante. Il faut rapporter :
\begin{enumerate}[leftmargin=*]
  \item erreur coefficient/opérateur en Frobenius et norme opérateur ;
  \item erreur solution $L^2$ et énergie $H^1$ ;
  \item résidu relatif et erreur par produit $Hv$ ;
  \item spectre de $B_\theta^{1/2}HB_\theta^{1/2}$ et conditionnement effectif ;
  \item stabilité couche par couche, facteur transitoire et extrapolation en
  nombre de boucles ;
  \item coût total, mémoire, paramètres, temps de préconditionnement et temps
  jusqu'à une tolérance ;
  \item calibration d'incertitude si l'inverse statistique est revendiqué.
\end{enumerate}

\subsection{Tests mécanistiques}

La trajectoire cachée doit être projetée vers les itérés Richardson, HeavyBall,
Chebyshev, CG et Newton sur exactement les mêmes prompts. Une corrélation
linéaire seule ne suffit pas : comparer aussi les résidus, les polynômes
spectraux par mode, les mises à jour et les contrefactuels obtenus en modifiant
le spectre sans changer les vecteurs propres. Des ablations doivent retirer la
mémoire $z_{\ell-1}$, partager ou non les poids, figer ou apprendre
$B_\theta$, et changer la distribution des seconds membres.

\section{Conclusion}

La littérature soutient trois constats robustes. Premièrement, les Transformers
peuvent apprendre ou représenter des algorithmes de régression en contexte, du
gradient préconditionné à des méthodes plus rapides. Deuxièmement, les
Transformers et neural operators sont devenus des surrogates puissants pour les
EDP et peuvent parfois s'adapter sans fine-tuning. Troisièmement, les approches
hybrides réseau--Krylov sont aujourd'hui les concurrentes naturelles lorsque la
précision numérique finale compte.

La contribution non couverte n'est donc pas « utiliser un Transformer pour une
EDP » ni « montrer qu'il peut faire Richardson ». Elle est de construire et
analyser un \emph{solveur-décodeur en contexte accéléré} : l'opérateur caché et
son préconditionnement sont appris à partir du prompt, la dynamique de
résolution est explicite et bouclée, les erreurs statistiques et numériques
sont séparées, et la théorie haute dimension prédit comment la structure de
faible rang apprise se transforme en convergence du solveur. HeavyBall fournit
le premier mécanisme pur le plus simple ; Chebyshev est le contrôle polynomial
nécessaire ; PCG reste la baseline numérique et le débouché hybride.

\clearpage
\begin{thebibliography}{99}
\small

\bibitem{brown2020}
T. B. Brown, B. Mann, N. Ryder, et al.
\newblock Language Models Are Few-Shot Learners.
\newblock \emph{NeurIPS}, 2020. \href{https://arxiv.org/abs/2005.14165}{arXiv:2005.14165}.

\bibitem{xie2022bayesian}
S. M. Xie, A. Raghunathan, P. Liang, and T. Ma.
\newblock An Explanation of In-Context Learning as Implicit Bayesian Inference.
\newblock \emph{ICLR}, 2022. \href{https://arxiv.org/abs/2111.02080}{arXiv:2111.02080}.

\bibitem{dai2023meta}
D. Dai, Y. Sun, L. Dong, et al.
\newblock Why Can GPT Learn In-Context? Language Models Implicitly Perform Gradient Descent as Meta-Optimizers.
\newblock 2023. \href{https://arxiv.org/abs/2212.10559}{arXiv:2212.10559}.

\bibitem{olsson2022}
C. Olsson, N. Elhage, N. Nanda, et al.
\newblock In-Context Learning and Induction Heads.
\newblock 2022. \href{https://arxiv.org/abs/2209.11895}{arXiv:2209.11895}.

\bibitem{chan2022distribution}
S. C. Y. Chan, A. Santoro, A. K. Lampinen, et al.
\newblock Data Distributional Properties Drive Emergent In-Context Learning in Transformers.
\newblock \emph{NeurIPS}, 2022. \href{https://arxiv.org/abs/2205.05055}{arXiv:2205.05055}.

\bibitem{chan2022weightscontext}
S. C. Y. Chan, I. Dasgupta, J. Kim, et al.
\newblock Transformers Generalize Differently from Information Stored in Context vs in Weights.
\newblock 2022. \href{https://arxiv.org/abs/2210.05675}{arXiv:2210.05675}.

\bibitem{min2022demonstrations}
S. Min, X. Lyu, A. Holtzman, et al.
\newblock Rethinking the Role of Demonstrations: What Makes In-Context Learning Work?
\newblock \emph{EMNLP}, 2022. \href{https://arxiv.org/abs/2202.12837}{arXiv:2202.12837}.

\bibitem{wies2023learnability}
N. Wies, Y. Levine, and A. Shashua.
\newblock The Learnability of In-Context Learning.
\newblock \emph{NeurIPS}, 2023. \href{https://arxiv.org/abs/2303.07895}{arXiv:2303.07895}.

\bibitem{garg2022}
S. Garg, D. Tsipras, P. Liang, and G. Valiant.
\newblock What Can Transformers Learn In-Context? A Case Study of Simple Function Classes.
\newblock \emph{NeurIPS}, 2022. \href{https://arxiv.org/abs/2208.01066}{arXiv:2208.01066}.

\bibitem{akyurek2023}
E. Akyürek, D. Schuurmans, J. Andreas, T. Ma, and D. Zhou.
\newblock What Learning Algorithm Is In-Context Learning? Investigations with Linear Models.
\newblock \emph{ICLR}, 2023. \href{https://arxiv.org/abs/2211.15661}{arXiv:2211.15661}.

\bibitem{vonoswald2023}
J. von Oswald, E. Niklasson, E. Randazzo, et al.
\newblock Transformers Learn In-Context by Gradient Descent.
\newblock \emph{ICML}, 2023. \href{https://arxiv.org/abs/2212.07677}{arXiv:2212.07677}.

\bibitem{ahn2023preconditioned}
K. Ahn, X. Cheng, H. Daneshmand, and S. Sra.
\newblock Transformers Learn to Implement Preconditioned Gradient Descent for In-Context Learning.
\newblock \emph{NeurIPS}, 2023. \href{https://arxiv.org/abs/2306.00297}{arXiv:2306.00297}.

\bibitem{mahankali2023}
A. Mahankali, T. B. Hashimoto, and T. Ma.
\newblock One Step of Gradient Descent Is Provably the Optimal In-Context Learner with One Layer of Linear Self-Attention.
\newblock 2023. \href{https://arxiv.org/abs/2307.03576}{arXiv:2307.03576}.

\bibitem{zhang2023trained}
R. Zhang, S. Frei, and P. L. Bartlett.
\newblock Trained Transformers Learn Linear Models In-Context.
\newblock 2023. \href{https://arxiv.org/abs/2306.09927}{arXiv:2306.09927}.

\bibitem{wu2024tasks}
J. Wu, D. Zou, Z. Chen, V. Braverman, Q. Gu, and P. L. Bartlett.
\newblock How Many Pretraining Tasks Are Needed for In-Context Learning of Linear Regression?
\newblock 2024. \href{https://arxiv.org/abs/2310.08391}{arXiv:2310.08391}.

\bibitem{raventos2023}
A. Raventós, M. Paul, F. Chen, and S. Ganguli.
\newblock Pretraining Task Diversity and the Emergence of Non-Bayesian In-Context Learning for Regression.
\newblock \emph{NeurIPS}, 2023. \href{https://arxiv.org/abs/2306.15063}{arXiv:2306.15063}.

\bibitem{li2023algorithms}
Y. Li, M. E. Ildiz, D. Papailiopoulos, and S. Oymak.
\newblock Transformers as Algorithms: Generalization and Stability in In-Context Learning.
\newblock \emph{ICML}, 2023. \href{https://arxiv.org/abs/2301.07067}{arXiv:2301.07067}.

\bibitem{bai2023statisticians}
Y. Bai, F. Chen, H. Wang, C. Xiong, and S. Mei.
\newblock Transformers as Statisticians: Provable In-Context Learning with In-Context Algorithm Selection.
\newblock \emph{NeurIPS}, 2023. \href{https://arxiv.org/abs/2306.04637}{arXiv:2306.04637}.

\bibitem{fu2024secondorder}
D. Fu, T.-Q. Chen, R. Jia, and V. Sharan.
\newblock Transformers Learn to Achieve Second-Order Convergence Rates for In-Context Linear Regression.
\newblock 2024. \href{https://arxiv.org/abs/2310.17086}{arXiv:2310.17086}.

\bibitem{vladymyrov2024}
M. Vladymyrov, J. von Oswald, M. Sandler, and R. Ge.
\newblock Linear Transformers Are Versatile In-Context Learners.
\newblock 2024. \href{https://arxiv.org/abs/2402.14180}{arXiv:2402.14180}.

\bibitem{cheng2024functional}
X. Cheng, Y. Chen, and S. Sra.
\newblock Transformers Implement Functional Gradient Descent to Learn Non-Linear Functions In Context.
\newblock 2024. \href{https://arxiv.org/abs/2312.06528}{arXiv:2312.06528}.

\bibitem{gatmiry2024learn}
K. Gatmiry, N. Saunshi, S. J. Reddi, S. Jegelka, and S. Kumar.
\newblock Can Looped Transformers Learn to Implement Multi-step Gradient Descent for In-context Learning?
\newblock 2024. \href{https://arxiv.org/abs/2410.08292}{arXiv:2410.08292}.

\bibitem{gatmiry2024depth}
K. Gatmiry, N. Saunshi, S. J. Reddi, S. Jegelka, and S. Kumar.
\newblock On the Role of Depth and Looping for In-Context Learning with Task Diversity.
\newblock 2024. \href{https://arxiv.org/abs/2410.21698}{arXiv:2410.21698}.

\bibitem{chen2024looped}
B. Chen, X. Li, Y. Liang, and Z. Shi.
\newblock Bypassing the Exponential Dependency: Looped Transformers Efficiently Learn In-Context by Multi-step Gradient Descent.
\newblock 2024. \href{https://arxiv.org/abs/2410.11268}{arXiv:2410.11268}.

\bibitem{huang2025cot}
J. Huang, Z. Wang, and J. D. Lee.
\newblock Transformers Learn to Implement Multi-step Gradient Descent with Chain of Thought.
\newblock 2025. \href{https://arxiv.org/abs/2502.21212}{arXiv:2502.21212}.

\bibitem{chen2024sparse}
X. Chen, L. Zhao, and D. Zou.
\newblock How Transformers Utilize Multi-Head Attention in In-Context Learning? A Case Study on Sparse Linear Regression.
\newblock 2024. \href{https://arxiv.org/abs/2408.04532}{arXiv:2408.04532}.

\bibitem{jin2024mixture}
Y. Jin, K. Balasubramanian, and L. Lai.
\newblock In-context Learning for Mixture of Linear Regressions: Existence, Generalization and Training Dynamics.
\newblock 2024. \href{https://arxiv.org/abs/2410.14183}{arXiv:2410.14183}.

\bibitem{guo2023representations}
T. Guo, W. Hu, S. Mei, et al.
\newblock How Do Transformers Learn In-Context Beyond Simple Functions? A Case Study on Learning with Representations.
\newblock 2023. \href{https://arxiv.org/abs/2310.10616}{arXiv:2310.10616}.

\bibitem{kwon2025subspace}
S. M. Kwon, A. S. Xu, C. Yaras, et al.
\newblock Out-of-Distribution Generalization of In-Context Learning: A Low-Dimensional Subspace Perspective.
\newblock 2025. \href{https://arxiv.org/abs/2505.14808}{arXiv:2505.14808}.

\bibitem{lu2025asymptotic}
Y. M. Lu, M. I. Letey, J. A. Zavatone-Veth, A. Maiti, and C. Pehlevan.
\newblock Asymptotic Theory of In-Context Learning by Linear Attention.
\newblock \emph{PNAS}, 2025. \href{https://arxiv.org/abs/2405.11751}{arXiv:2405.11751}.

\bibitem{letey2026alignment}
M. I. Letey, J. A. Zavatone-Veth, Y. M. Lu, and C. Pehlevan.
\newblock Pretrain-Test Task Alignment Governs Generalization in In-Context Learning.
\newblock 2025/2026. \href{https://arxiv.org/abs/2509.26551}{arXiv:2509.26551}.

\bibitem{takanami2026}
K. Takanami, T. Takahashi, and Y. Kabashima.
\newblock Learning Linear Regression with Low-Rank Tasks In-Context.
\newblock \emph{AISTATS}, 2026. \href{https://arxiv.org/abs/2510.04548}{arXiv:2510.04548}.

\bibitem{yan2026richardson}
M. Yan, D. Li, C. Kulick, and S. Tang.
\newblock Transformers Can Implement Preconditioned Richardson Iteration for In-Context Gaussian Kernel Regression.
\newblock 2026. \href{https://arxiv.org/abs/2605.08475}{arXiv:2605.08475}.

\bibitem{cole2025systems}
F. Cole, Y. Lu, W. Xu, and T. Zhang.
\newblock In-Context Learning of Linear Systems: Generalization Theory and Applications to Operator Learning.
\newblock 2024/2025. \href{https://arxiv.org/abs/2409.12293}{arXiv:2409.12293}.

\bibitem{cole2025dynamics}
F. Cole, Y. Zhao, Y. Lu, and T. Zhang.
\newblock In-Context Learning of Linear Dynamical Systems with Transformers: Approximation Bounds and Depth-Separation.
\newblock 2025. \href{https://arxiv.org/abs/2502.08136}{arXiv:2502.08136}.

\bibitem{goel2024kalman}
G. Goel and P. Bartlett.
\newblock Can a Transformer Represent a Kalman Filter?
\newblock \emph{L4DC}, 2024. \href{https://proceedings.mlr.press/v242/goel24a.html}{PMLR 242}.

\bibitem{akram2024state}
U. Akram and H. Vikalo.
\newblock Transformers as Implicit State Estimators: In-Context Learning in Dynamical Systems.
\newblock 2024. \href{https://arxiv.org/abs/2410.16546}{arXiv:2410.16546}.

\bibitem{lu2021deeponet}
L. Lu, P. Jin, G. Pang, Z. Zhang, and G. E. Karniadakis.
\newblock Learning Nonlinear Operators via DeepONet Based on the Universal Approximation Theorem of Operators.
\newblock \emph{Nature Machine Intelligence}, 2021. \href{https://arxiv.org/abs/1910.03193}{arXiv:1910.03193}.

\bibitem{li2021fno}
Z. Li, N. Kovachki, K. Azizzadenesheli, et al.
\newblock Fourier Neural Operator for Parametric Partial Differential Equations.
\newblock \emph{ICLR}, 2021. \href{https://arxiv.org/abs/2010.08895}{arXiv:2010.08895}.

\bibitem{kovachki2023neuraloperator}
N. Kovachki, Z. Li, B. Liu, et al.
\newblock Neural Operator: Learning Maps Between Function Spaces with Applications to PDEs.
\newblock \emph{JMLR}, 2023. \href{https://arxiv.org/abs/2108.08481}{arXiv:2108.08481}.

\bibitem{lanthaler2022error}
S. Lanthaler, S. Mishra, and G. E. Karniadakis.
\newblock Error Estimates for DeepONets: A Deep Learning Framework in Infinite Dimensions.
\newblock \emph{Transactions of Mathematics and Its Applications}, 2022. \href{https://arxiv.org/abs/2102.09618}{arXiv:2102.09618}.

\bibitem{kim2022rademacher}
T. Kim and M. Kang.
\newblock Bounding the Rademacher Complexity of Fourier Neural Operators.
\newblock 2022. \href{https://arxiv.org/abs/2209.05150}{arXiv:2209.05150}.

\bibitem{cao2021galerkin}
S. Cao.
\newblock Choose a Transformer: Fourier or Galerkin.
\newblock \emph{NeurIPS}, 2021. \href{https://arxiv.org/abs/2105.14995}{arXiv:2105.14995}.

\bibitem{li2022oformer}
Z. Li, K. Meidani, and A. B. Farimani.
\newblock Transformer for Partial Differential Equations' Operator Learning.
\newblock 2022. \href{https://arxiv.org/abs/2205.13671}{arXiv:2205.13671}.

\bibitem{hao2023gnot}
Z. Hao, Z. Wang, H. Su, et al.
\newblock GNOT: A General Neural Operator Transformer for Operator Learning.
\newblock \emph{ICML}, 2023. \href{https://arxiv.org/abs/2302.14376}{arXiv:2302.14376}.

\bibitem{wu2024transolver}
H. Wu, H. Luo, H. Wang, et al.
\newblock Transolver: A Fast Transformer Solver for PDEs on General Geometries.
\newblock \emph{ICML}, 2024. \href{https://arxiv.org/abs/2402.02366}{arXiv:2402.02366}.

\bibitem{hao2024dpot}
Z. Hao, C. Su, S. Liu, et al.
\newblock DPOT: Auto-Regressive Denoising Operator Transformer for Large-Scale PDE Pre-Training.
\newblock 2024. \href{https://arxiv.org/abs/2403.03542}{arXiv:2403.03542}.

\bibitem{herde2024poseidon}
M. Herde, B. Raonić, T. Rohner, et al.
\newblock Poseidon: Efficient Foundation Models for PDEs.
\newblock 2024. \href{https://arxiv.org/abs/2405.19101}{arXiv:2405.19101}.

\bibitem{ye2024pdeformer}
Z. Ye, X. Huang, L. Chen, et al.
\newblock PDEformer: Towards a Foundation Model for One-Dimensional Partial Differential Equations.
\newblock 2024. \href{https://arxiv.org/abs/2402.12652}{arXiv:2402.12652}.

\bibitem{shen2024ups}
J. Shen, T. Marwah, and A. Talwalkar.
\newblock UPS: Efficiently Building Foundation Models for PDE Solving via Cross-Modal Adaptation.
\newblock 2024. \href{https://arxiv.org/abs/2403.07187}{arXiv:2403.07187}.

\bibitem{lu2024unisolver}
H. Zhou, Y. Ma, H. Wu, H. Wang, and M. Long.
\newblock Unisolver: PDE-Conditional Transformers Towards Universal Neural PDE Solvers.
\newblock 2024. \href{https://arxiv.org/abs/2405.17527}{arXiv:2405.17527}.

\bibitem{yang2023icon}
L. Yang, S. Liu, T. Meng, and S. J. Osher.
\newblock In-Context Operator Learning with Data Prompts for Differential Equation Problems.
\newblock \emph{PNAS}, 2023/2024. \href{https://arxiv.org/abs/2304.07993}{arXiv:2304.07993}.

\bibitem{yang2024icongeneralization}
L. Yang and S. J. Osher.
\newblock PDE Generalization of In-Context Operator Networks: A Study on 1D Scalar Nonlinear Conservation Laws.
\newblock 2024. \href{https://arxiv.org/abs/2401.07364}{arXiv:2401.07364}.

\bibitem{mishra2025continuum}
A. Mishra, Y. Patel, and A. Tewari.
\newblock Continuum Transformers Perform In-Context Learning by Operator Gradient Descent.
\newblock 2025. \href{https://arxiv.org/abs/2505.17838}{arXiv:2505.17838}.

\bibitem{ovadia2023vito}
O. Ovadia, A. Kahana, P. Stinis, E. Turkel, and G. E. Karniadakis.
\newblock ViTO: Vision Transformer-Operator.
\newblock 2023. \href{https://arxiv.org/abs/2303.08891}{arXiv:2303.08891}.

\bibitem{molinaro2023nio}
R. Molinaro, Y. Yang, B. Engquist, and S. Mishra.
\newblock Neural Inverse Operators for Solving PDE Inverse Problems.
\newblock 2023. \href{https://arxiv.org/abs/2301.11167}{arXiv:2301.11167}.

\bibitem{cao2024bilo}
R. Z. Zhang, C. E. Miles, X. Xie, and J. S. Lowengrub.
\newblock BiLO: Bilevel Local Operator Learning for PDE Inverse Problems.
\newblock 2024. \href{https://arxiv.org/abs/2404.17789}{arXiv:2404.17789}.

\bibitem{wang2024lno}
T. Wang and C. Wang.
\newblock Latent Neural Operator for Solving Forward and Inverse PDE Problems.
\newblock 2024. \href{https://arxiv.org/abs/2406.03923}{arXiv:2406.03923}.

\bibitem{sappl2019preconditioners}
J. Sappl, L. Seiler, M. Harders, and W. Rauch.
\newblock Deep Learning of Preconditioners for Conjugate Gradient Solvers in Urban Water Related Problems.
\newblock 2019. \href{https://arxiv.org/abs/1906.06925}{arXiv:1906.06925}.

\bibitem{li2023preconditioner}
Y. Li, P. Y. Chen, T. Du, and W. Matusik.
\newblock Learning Preconditioner for Conjugate Gradient PDE Solvers.
\newblock 2023. \href{https://arxiv.org/abs/2305.16432}{arXiv:2305.16432}.

\bibitem{hausner2023neuralif}
P. Häusner, O. Öktem, and J. Sjölund.
\newblock Neural Incomplete Factorization: Learning Preconditioners for the Conjugate Gradient Method.
\newblock 2023. \href{https://arxiv.org/abs/2305.16368}{arXiv:2305.16368}.

\bibitem{rudikov2024fcgno}
A. Rudikov, V. Fanaskov, E. Muravleva, Y. Laevsky, and I. Oseledets.
\newblock Neural Operators Meet Conjugate Gradients: The FCG-NO Method for Efficient PDE Solving.
\newblock 2024. \href{https://arxiv.org/abs/2402.05598}{arXiv:2402.05598}.

\bibitem{han2024ugrid}
X. Han, F. Hou, and H. Qin.
\newblock UGrid: An Efficient-and-Rigorous Neural Multigrid Solver for Linear PDEs.
\newblock 2024. \href{https://arxiv.org/abs/2408.04846}{arXiv:2408.04846}.

\bibitem{greenfeld2019multigrid}
D. Greenfeld, M. Galun, R. Basri, I. Yavneh, and R. Kimmel.
\newblock Learning to Optimize Multigrid PDE Solvers.
\newblock \emph{ICML}, 2019. \href{https://proceedings.mlr.press/v97/greenfeld19a.html}{PMLR 97}.

\bibitem{engl1996inverse}
H. W. Engl, M. Hanke, and A. Neubauer.
\newblock \emph{Regularization of Inverse Problems}.
\newblock Kluwer, 1996.

\bibitem{stuart2010bayesian}
A. M. Stuart.
\newblock Inverse Problems: A Bayesian Perspective.
\newblock \emph{Acta Numerica}, 2010.

\bibitem{polyak1964}
B. T. Polyak.
\newblock Some Methods of Speeding Up the Convergence of Iteration Methods.
\newblock \emph{USSR Computational Mathematics and Mathematical Physics}, 1964.

\bibitem{saad2003}
Y. Saad.
\newblock \emph{Iterative Methods for Sparse Linear Systems}, second edition.
\newblock SIAM, 2003.

\bibitem{barbier2019adaptive}
J. Barbier and N. Macris.
\newblock The Adaptive Interpolation Method: A Simple Scheme to Prove Replica Formulas in Bayesian Inference.
\newblock \emph{Probability Theory and Related Fields}, 2019.

\bibitem{bayati2011amp}
M. Bayati and A. Montanari.
\newblock The Dynamics of Message Passing on Dense Graphs, with Applications to Compressed Sensing.
\newblock \emph{IEEE Transactions on Information Theory}, 2011.

\end{thebibliography}

\end{document}
